\documentclass[11pt]{article}
\usepackage[margin=1in]{geometry}
\usepackage{amsmath,amssymb,amsthm,hyperref}
\newtheorem{theorem}{Theorem}
\newtheorem{lemma}{Lemma}
\title{A quadratic counterexample to the proposed real B\'ezout bound}
\author{ziangni-sys}
\date{Conjecture 00000002333}
\begin{document}
\sloppy
\maketitle
\noindent The English and Chinese statements define the real B\'ezout number as
an upper bound for the number of solutions of real algebraic systems and propose
\[
 B(d,m)=\left(\frac{d}{\sqrt m}\right)^m e^{1/2}.
\]
Here $m$ is the number of variables (and equations in the square system below),
and $d$ is the maximal total degree. We disprove this explicit upper bound.
The second assertion, about tightness for balanced random systems, cannot
restore a bound that already fails on a real square system of equal degrees.
No probability distribution on random systems is assumed or needed.

\begin{theorem}
The real polynomial system
\[
 f_0(x_0,x_1)=x_0^2-1=0,\qquad
 f_1(x_0,x_1)=x_1^2-1=0
\]
has precisely four isolated real solutions. Both equations have total degree
$2$, whereas $B(2,2)<4$.
\end{theorem}
\begin{proof}
For each $i\in\{0,1\}$, the equation $x_i^2-1=0$ factors as
$(x_i-1)(x_i+1)=0$. Over $\mathbb R$ this is equivalent to $x_i=1$ or $x_i=-1$.
Consequently the complete solution set is
\[
 \{(-1,-1),(-1,1),(1,-1),(1,1)\}.
\]
These four vectors are distinct. Each polynomial contains the nonzero monomial
$x_i^2$ and a constant, so its total degree is exactly $2$.

If two distinct solutions differ, they differ in some coordinate by $2$.
Their distance in the supremum norm is therefore $2$. In particular, the
open ball of radius $1$ about any solution contains no other solution. Thus
all four solutions are isolated; this example also meets the customary
restriction of B\'ezout bounds to isolated roots.

Since $(\sqrt2)^2=2$,
\[
 B(2,2)=\frac{4}{2}e^{1/2}=2e^{1/2}.
\]
The standard exponential identity gives $(e^{1/2})^2=e$. The rigorous bound
$e<2.7182818286<4$, together with $e^{1/2}>0$, implies $e^{1/2}<2$.
Hence $B(2,2)<4$, strictly below the actual number of isolated real roots.
\end{proof}

\newpage
\section*{Formal verification and correspondence}
The Lean project uses genuine multivariate polynomials over $\mathbb R$, with
variable type \texttt{Fin 2}. Its definition \texttt{equation i} is
\texttt{MvPolynomial.X i \^{} 2} minus \texttt{MvPolynomial.C 1}.
The theorem \texttt{equation\_degree} proves that each polynomial's actual
\texttt{totalDegree} is $2$, and \texttt{equation\_eval} computes its actual
polynomial evaluation.

The predicate \texttt{IsSolution} requires all these polynomial evaluations
to vanish. The theorem \texttt{solution\_iff} proves the complete coordinate
classification. The equivalence \texttt{rootEquiv} identifies the actual
solution subtype with the Boolean functions \texttt{Fin 2} $\to$ \texttt{Bool}:
a Boolean choice specifies each sign. Both inverse laws are proved.
\texttt{solution\_count} therefore computes the cardinality of the entire
solution set as $4$, rather than counting an asserted list of roots.

The theorem \texttt{roots\_isolated} states that two solution vectors at
supremum-norm distance less than $1$ must be equal. It proves this by bounding
each coordinate difference by the actual norm of the vector difference,
and then using the full coordinate classification. This is an explicit
isolation radius for every root in the formalized metric space.

The definition \texttt{proposedBound} is the displayed source formula with
natural parameters cast to real numbers and the actual real square root and
exponential. \texttt{bound\_at\_two} proves its exact specialization.
\texttt{exp\_half\_lt\_two} derives the strict exponential inequality from
Mathlib's certified bound \texttt{Real.exp\_one\_lt\_d9}.
The final theorem \texttt{claimed\_bound\_fails} proves
\[
 B(2,2)<\#\{x\in\mathbb R^2:f_0(x)=f_1(x)=0\}.
\]

\section*{Scope}
This is a disproof of the explicit universal upper-bound clause, with
positive degree and dimension and finitely many isolated real roots.
It does not replace that clause by an expected root count under an
unspecified random model. Both versions of the original statement explicitly
call the quantity an upper bound for the solution count. The two equations
have the same degree, so unequal-degree conventions cause no issue.
No restriction excluding these polynomials appears in either source version.

\section*{Reproduction}
The project pins Lean 4.19.0 and Mathlib commit
\begin{center}\small\texttt{c44e0c8ee63ca166450922a373c7409c5d26b00b}\end{center}
From \texttt{lean/}, run \texttt{lake update}, \texttt{lake exe cache get},
then \texttt{lake build}. Check the source directly with
\begin{center}\small\texttt{lake env lean -DwarningAsError=true Main.lean}.\end{center}
The complete proof contains no admitted statements or additional axioms.
The accompanying verification logs record compilation and the axiom audit.
\end{document}
